\section*{Chapitre \ref{ch:ml:sequence-models-on-order-books} --- Modèles de séquences sur les carnets d'ordres}

\begin{solution}{exo:ml:sequence-models-on-order-books:1}
$1 + (3 - 1)(1 + 2 + 4 + 8) = 31$ instantanés~; avec une largeur de cinq, $1 + 4\times15 = 61$.
\end{solution}

\begin{solution}{exo:ml:sequence-models-on-order-books:2}
Toujours «~neutre~»~: 72.7~\%. En tirant au hasard selon les proportions des classes~: $0.727^2 + 0.147^2 + 0.125^2 = 56.6~\%$.
\end{solution}

\begin{solution}{exo:ml:sequence-models-on-order-books:3}
$T\times T\times h = 20\times20\times32 = 12\,800$ multiplications-additions~; quadratique en la longueur de la séquence (et linéaire en la
largeur).
\end{solution}

\begin{solution}{exo:ml:sequence-models-on-order-books:4}
Moyenner les sorties de la convolution sur la fenêtre revient à pondérer l'instantané d'il y a vingt secondes autant que l'instantané courant,
alors que seul l'état récent prédit les secondes suivantes. Mieux vaut lire à la place la dernière position (comme le TCN et le \omterm{def:ml:sequence-models-on-order-books:attention}{transformer}
le font), ou aplatir l'axe temporel dans le classifieur.
\end{solution}

\begin{solution}{exo:ml:sequence-models-on-order-books:5}
La moyenne des prix milieu des dix secondes suivantes évolue moins et plus régulièrement que le mouvement suivant, et ses premières secondes prolongent le
flux courant, que voient les entrées~; une grande part de sa prévisibilité vient de ce recouvrement. Une transaction rapporte la variation de prix jusqu'à sa sortie
et non une moyenne des prix milieu futurs, si bien qu'un modèle évalué sur l'étiquette lissée surestime l'edge exploitable~; il faut l'évaluer
sur le P\&L de la transaction qu'il piloterait (chapitre~29).
\end{solution}

\begin{solution}{exo:ml:sequence-models-on-order-books:6}
Des découpages aléatoires placent dans les deux ensembles des fenêtres d'un même jour qui se chevauchent dans le temps~: le réseau est testé sur des quasi-copies de
ses fenêtres d'entraînement (chapitre~3). Or 88~\% se compare à la prédiction systématique de la classe majoritaire. Découper par jour,
purger les bords, et rapporter un IC ou un P\&L.
\end{solution}

\begin{solution}{exo:ml:sequence-models-on-order-books:7}
\texttt{scaling(16, name='LSTM')}~: IC 0.381, contre 0.441 pour le TCN et 0.568 pour la \omterm{def:ml:linear-and-regularised-baselines:logistic}{régression logistique} sur les mêmes
seize sessions.
\end{solution}

\begin{solution}{exo:ml:sequence-models-on-order-books:8}
Une couche causale de dilatation $r$ à la position $t$ lit $t, t - r,\dots, t - (k-1)r$~: elle étend la portée vers le passé de
$(k-1)r$, et jamais vers l'avant. Empiler des couches additionne les portées~: $(k-1)(1 + 2 + \dots + 2^{L-1}) = (k-1)(2^L - 1)$, plus la
position elle-même. Par récurrence sur les couches, toute sortie en $t$ est une fonction des entrées en $t$ ou antérieures~; les connexions résiduelles
ajoutent que la même position.
\end{solution}

\begin{solution}{pb:ml:sequence-models-on-order-books:1}
\begin{enumerate}
\item Un instantané~: pour chacun des dix niveaux de chaque côté, la distance au prix milieu en ticks et la taille logarithmique en lots (40 nombres)~;
une fenêtre~: les 20 derniers instantanés, soit dix secondes.
\item Moyenne du prix milieu sur les dix secondes suivantes moins le prix milieu courant, neutre à un quart de tick près. Parts de l'\omterm{def:ml:why-financial-machine-learning-is-different:sets}{ensemble de test}~: 72.7~\% neutre,
14.7~\% en baisse, 12.5~\% en hausse.
\item Prédire systématiquement «~neutre~» a raison dans 72.7~\% des cas~; le chapitre classe les modèles par l'IC de $\P(\text{hausse}) -
\P(\text{baisse})$ avec la variation future.
\item Les fenêtres d'une même session se recouvrent et partagent son état~: d'un découpage à l'autre, elles fuient.
\item CNN~: neuf instantanés par pile de filtres, puis une moyenne sur les 20~; LSTM et \omterm{def:ml:sequence-models-on-order-books:attention}{transformer} les 20~; TCN 31, plus que la
fenêtre.
\item CNN 0.727 et 0.060~; LSTM 0.729 et 0.292~; TCN 0.766 et 0.366~; attention 0.775 et 0.426.
\item Exactitude 0.781, IC 0.543.
\item Écarts-types de 0.15 à 0.20 pour l'IC des modèles de séquences, 0.09 pour la \omterm{def:ml:linear-and-regularised-baselines:logistic}{régression logistique}.
\item Après une à quatre \omterm{def:ml:neural-networks-for-noisy-tabular-data:backprop}{époques}~: 8\,808 fenêtres sont peu pour des modèles qui lisent 800 nombres chacun.
\item 0.13, 0.30 et 0.44.
\item À peine (0.54 à 0.57)~: ses dix variables résument déjà ce que le chapitre dit des secondes suivantes~; il a peu de
paramètres, et peu de variance à réduire.
\item Continuer de monter au-delà de 0.57 à mesure que des sessions s'ajoutent, ce que le doublement de huit à seize ne montre pas encore.
\item Que des réseaux plus simples nourris de variables de flux d'ordres surpassent la plupart des modèles entraînés directement sur les carnets d'ordres.
\item Que les quinze modèles publiés perdent une grande part de leur performance sur des données nouvelles.
\item Des milliards de cotations regroupées sur de nombreuses actions.
\item \emph{Résultat nommé}~: le meilleur modèle de séquences (attention) atteint un IC de 0.426 et une exactitude de 0.775, contre
0.543 et 0.781 pour la \omterm{def:ml:linear-and-regularised-baselines:logistic}{régression logistique}, avec des écarts-types entre sessions et graines de 0.15 contre 0.09~: un gain de $-0.12$
d'IC, l'écart se resserrant à mesure que des sessions s'ajoutent.
\item La \omterm{def:ml:linear-and-regularised-baselines:logistic}{régression logistique} (ou les arbres boostés) sur les variables de flux d'ordres~; ensuite, le TCN avec les mêmes variables en canaux
supplémentaires, entraîné sur bien plus de sessions et d'instruments.
\item Grâce au chapitre~2~: un seuil méta-étiqueté sur $\P(\text{hausse}) - \P(\text{baisse})$ qui ne paie le spread que lorsque le
mouvement attendu le dépasse, évalué sur le P\&L des transactions.
\item Davantage de sessions et d'instruments à regrouper, et des entrées market-by-order~; les mêmes règles de découpage.
\item Bien plus de données que la variable n'en exige, ou la variable elle-même en entrée.
\end{enumerate}
\end{solution}

\begin{solution}{iq:ml:sequence-models-on-order-books:1}
Une convolution dont la sortie en $t$ n'utilise que les entrées en $t$ et avant. Un modèle de trading ne doit pas voir le futur de sa propre
fenêtre~; une convolution centrée, entraînée hors ligne, fait fuiter les pas suivants dans chaque prédiction.
\iqlookfor{la définition et le biais d'anticipation qu'elle empêche.}
\end{solution}

\begin{solution}{iq:ml:sequence-models-on-order-books:2}
$\mathsf Q = ZA_{\mathsf Q}$, $\mathsf K = ZA_{\mathsf K}$, $\mathsf V = ZA_{\mathsf V}$, avec pour sortie $\mathrm{softmax}(\mathsf
Q\mathsf K^\top/\sqrt h)\mathsf V$~: chaque position moyenne les valeurs de toutes les positions, avec des poids issus de la similarité requête--clé.
Coût $O(T^2h)$ en temps et $O(T^2)$ en mémoire pour les poids.
\iqlookfor{la formule, la mise à l'échelle et le coût quadratique.}
\end{solution}

\begin{solution}{iq:ml:sequence-models-on-order-books:3}
Quelle est la part de la classe majoritaire~; comment l'étiquette est construite (lissage, horizon, bande morte)~; comment les données ont été découpées (par jour,
avec purge)~; résultats par jour et par instrument, avec les fourchettes~; un IC ou un P\&L après coûts~; comparaison avec un modèle linéaire sur
variables de flux d'ordres.
\iqlookfor{équilibre des classes, construction de l'étiquette, découpage et \omterm{def:ml:linear-and-regularised-baselines:baseline}{modèle de référence}.}
\end{solution}

\begin{solution}{iq:ml:sequence-models-on-order-books:4}
Les éléments disponibles (Kolm, Turiel et Westray~; la courbe d'apprentissage de ce chapitre) privilégient les variables de flux d'ordres à des tailles de données réalistes~:
elles sont stationnaires et résument déjà le carnet d'ordres. Les niveaux bruts peuvent l'emporter avec de très grands jeux de données regroupés (Sirignano et
Cont). Commencer par les variables~; ajouter les entrées brutes en canaux supplémentaires.
\iqlookfor{la taille des données comme facteur déterminant, avec les résultats publiés.}
\end{solution}

\begin{solution}{iq:ml:sequence-models-on-order-books:5}
Par sessions entières (jours), dans l'ordre chronologique, la période de test venant après l'entraînement et la validation, purgée aux bords~; jamais
fenêtre par fenêtre au hasard~; la normalisation est ajustée sur les jours d'entraînement.
\iqlookfor{découpages au niveau des sessions et raisons de la fuite des fenêtres.}
\end{solution}

\begin{solution}{iq:ml:sequence-models-on-order-books:6}
Sur des fenêtres courtes, les instantanés les plus récents portent le signal~; un TCN les lit directement et s'entraîne en parallèle, alors qu'un
LSTM doit les porter à travers son état et est plus lent à ajuster. Un LSTM peut faire mieux lorsque l'historique pertinent est long et
irrégulier, avec un état à accumuler (une file qui se constitue lentement, un métaordre).
\iqlookfor{récence, dynamique d'entraînement, et moment où la récurrence est rentable.}
\end{solution}
